\section{KPI Definitions}
\label{app:kpis}

All definitions below are as implemented in \texttt{scripts/eval\_timetable\_kpis.py}. Let $A$ be
the set of assigned classes in the solution, $S$ the set of students in the enrolment data, and
for a student $s$ let $E(s)$ be the set of assigned classes $s$ is enrolled in.

\subsection{The overlap predicate}

Two classes $c_1, c_2$ overlap, written $\mathrm{ov}(c_1,c_2)$, if and only if all three hold:

\begin{enumerate}
  \item they share at least one day: $\mathrm{days}(c_1) \land \mathrm{days}(c_2) \neq 0$, where
    days are a seven-bit mask with index 0 = Monday;
  \item their slot intervals intersect:
    $\mathrm{start}(c_1) < \mathrm{end}(c_2) \;\wedge\; \mathrm{start}(c_2) < \mathrm{end}(c_1)$,
    so intervals that merely touch at an endpoint do \emph{not} overlap;
  \item if both carry date patterns, the patterns share at least one active week.
\end{enumerate}

Room and instructor play no part in the test. Slots are five minutes when the session declares 288
slots per day.

\subsection{Hard checks}

Each check reports \emph{checked}, \emph{passed}, \emph{failed} and a pass rate, together with up
to eight example failures.

\begin{table}[H]
\centering
\caption{The ten hard checks.}
\label{tab:hardkpi}
\footnotesize
\begin{tabular}{@{}L{3.5cm}L{11.4cm}@{}}
\toprule
\textbf{Check} & \textbf{Definition} \\
\midrule
\texttt{coverage} &
  Fraction of classes declared in the offering file that carry an assignment in the solution. \\
\texttt{time\_pattern\_fidelity} &
  Fraction of assigned classes whose number of meetings and minutes per meeting match the time
  pattern declared for their subpart. \\
\texttt{date\_pattern} &
  Fraction of assigned classes whose date pattern matches the one declared in the preference file. \\
\texttt{room\_capacity} &
  Fraction of assigned classes whose enrolment limit does not exceed the capacity of the assigned
  room. \\
\texttt{lunch\_weekend} &
  Fraction of assigned classes that do not fall on a weekend or inside a blocked slot. \\
\texttt{room\_unavailable} &
  Fraction of assigned classes placed in a room that is available at that time (XML mode). \\
\texttt{instructor\_lunch\_breaks} &
  Fraction of instructors who retain at least one free 60-minute window within 11:30--15:30 on
  every teaching day. \\
\texttt{student\_lunch\_breaks} &
  The same test applied to each student's combined timetable. \\
\texttt{instructor\_double\_booking} &
  Fraction of classes not overlapping another class with the same instructor. \\
\texttt{room\_double\_booking} &
  Fraction of classes not overlapping another class in the same room. \\
\texttt{student\_conflicts} &
  Fraction of students $s \in S$ for which no pair $c_1,c_2 \in E(s)$ satisfies
  $\mathrm{ov}(c_1,c_2)$. \\
\bottomrule
\end{tabular}
\end{table}

\subsection{Student conflict quantities}

Four related quantities are reported and must not be interchanged:

\begin{align*}
  \text{students\_with\_conflict} &= \bigl|\{\, s \in S \;:\; \exists\, c_1,c_2 \in E(s),\;
      \mathrm{ov}(c_1,c_2) \,\}\bigr| \\[2pt]
  \text{student\_conflict\_pairs} &= \bigl|\{\, (s,c_1,c_2) \;:\; c_1,c_2 \in E(s),\;
      \mathrm{ov}(c_1,c_2) \,\}\bigr| \\[2pt]
  C &= \frac{\text{students\_with\_conflict}}{\bigl|\{\, s : |E(s)| \ge 2 \,\}\bigr|}
      \qquad\text{(violation rate, students at risk)} \\[2pt]
  C_{\mathrm{checker}} &= \frac{\text{students\_with\_conflict}}{|S|}
      \qquad\text{(violation rate over all students)}
\end{align*}

$C$ and $C_{\mathrm{checker}}$ coincide whenever every student is enrolled in at least two classes,
which is true of the current dataset. Section~\ref{sec:experiments} reports $C$.

\subsection{Preference checks}

\begin{itemize}
  \item \texttt{room\_preference}: fraction of assigned classes placed in a room that appears in
    the class's preferred-room list.
  \item \texttt{instructor\_identity}: fraction of assigned classes whose instructor matches the
    one expected from the source data. \emph{This check is currently unreliable} because the two
    sides use different name formats (Section~\ref{sec:evaluation}).
  \item \texttt{unav\_onl}: count and percentage of assignments landing in the unavailable or
    online room pools.
\end{itemize}

\subsection{Aggregate scores}

Both aggregates are equal-weight means of the component pass rates, on a 0--100 scale:
\[
  \text{hard feasibility} = \frac{100}{|H|}\sum_{h \in H} r_h ,
  \qquad
  \text{preference score} = \frac{100}{|P|}\sum_{p \in P} r_p ,
\]
where $H$ is the set of hard checks and $P$ the set of preference checks, and $r$ is a pass rate in
$[0,1]$. Because the weights are equal, a single broken check moves the aggregate substantially;
the component rates should always be read alongside it.

\subsection{Efficiency metrics}

These are descriptive and carry no pass/fail threshold:

\begin{itemize}
  \item \textbf{Instructor weekly load}: mean, standard deviation, minimum and maximum of assigned
    minutes per week per instructor.
  \item \textbf{Travel violations}: consecutive classes for one student or instructor in buildings
    further apart than the gap between them permits.
  \item \textbf{Slot spread}: proportion of assignments in the early, middle and late thirds of the
    teaching day.
  \item \textbf{Back-to-back density}: proportion of assignments immediately adjacent to another
    class for the same instructor or cohort.
  \item \textbf{Unavailable/online usage}: share of assignments in those room pools.
\end{itemize}
